\begin{notechapter}
\chapter{The Flat Circle \texorpdfstring{$\R/\Z$}{R/Z}}\label{ch:flat-circle}
\chaptersubtitle{A circle, a group, and a curve with zero curvature}

The pitfall of \cref{gradient-tangent-space:sec:affine} uses $\R/\Z$ as a space that is flat but
not affine. This appendix explains the three claims hidden in the phrase ``the circle of length $1$,
$\R/\Z$, is flat'': that $\R/\Z$ is a circle, that its length is $1$, and that its curvature is zero
although it closes up on itself. It ends with what this means for parameters that are angles.

\section{The quotient}

Call $x, y\in\R$ equivalent when $y - x\in\Z$. This is an equivalence relation, and $\R/\Z$ is the set of
its classes $[x] = x + \Z = \{\dots, x - 1, x, x + 1, \dots\}$, with the projection $\pi:\R\to\R/\Z$, $\pi(x) = [x]$.
\begin{itemize}[nosep]
  \item \emph{Representatives.} Every class has exactly one representative in $[0, 1)$, namely
        $x - \lfloor x\rfloor$. As a set, $\R/\Z$ is ``$[0, 1)$'', but with $1$ identified with $0$.
  \item \emph{Group.} $[x] + [y] = [x + y]$ is well defined (if $x' - x$ and $y' - y$ are integers, so is their
        sum), which makes $\R/\Z$ an abelian group and $\pi$ a group homomorphism with kernel $\Z$.
  \item \emph{Topology.} A set $U\subset\R/\Z$ is declared open when $\pi^{-1}(U)$ is open in $\R$ (the quotient
        topology). Then $\R/\Z = \pi([0, 1])$ is compact, as the continuous image of a compact set.
\end{itemize}

\section{It is a circle}

\begin{proposition}\label{flat-circle:prop:circle}
Let $S^1 = \{z\in\mathbb C : \abs z = 1\}$ and $p(x) = e^{2\pi ix}$. Then $\bar p([x]) = p(x)$ defines a map
$\bar p:\R/\Z\to S^1$ that is both a group isomorphism from $(\R/\Z, +)$ to $(S^1, \times)$ and a homeomorphism.
\end{proposition}

\begin{proof}
$p(x) = p(y)\iff e^{2\pi i(y - x)} = 1\iff y - x\in\Z$. The implication from right to left says that $\bar p$ is well
defined, the one from left to right that it is injective. It is surjective because every unit
complex number is $e^{i\alpha}$ for some $\alpha$. It is a homomorphism because
$e^{2\pi i(x + y)} = e^{2\pi ix}e^{2\pi iy}$. It is continuous because $\bar p\circ\pi = p$ is continuous and $\R/\Z$ carries the
quotient topology. Finally, a continuous bijection from a compact space to a Hausdorff space is a
homeomorphism (it maps closed, hence compact, sets to compact, hence closed, sets).
\end{proof}

Topologically, $\R/\Z$ is the segment $[0, 1]$ with its two ends glued together
(\cref{flat-circle:fig:circle}); algebraically, adding in $\R/\Z$ is rotating the circle.

\begin{figure}[ht]
\centering
\begin{tikzpicture}[>=Stealth, scale=1.35]
  % the real line and a fundamental domain
  \draw[->, gray] (-1.3,0) -- (4.6,0);
  \foreach \x/\l in {0/0, 3/1}{\draw (\x,0.08) -- (\x,-0.08) node[below] {$\l$};}
  \draw[very thick, bookaccent] (0,0) -- (3,0);
  \fill[bookpitfall] (0,0) circle (2pt); \fill[bookpitfall] (3,0) circle (2pt);
  \fill[bookdef] (0.75,0) circle (2pt) node[above, black] {$\tfrac14$};
  \fill[bookexample] (1.5,0) circle (2pt) node[above, black] {$\tfrac12$};
  \fill[bookpitfall!40] (-1,0) circle (1.6pt) node[below, black] {$-1$};
  \fill[bookpitfall!40] (4,0) circle (1.6pt) node[below, black] {$2$};
  \node[gray] at (1.5,-0.75) {$\R$};
  \draw[->, thick] (5.0,0) -- node[above] {$c$} (6.0,0);
  % the circle of length 1 (radius 1/(2 pi), drawn at the scale of the segment)
  \begin{scope}[shift={(7.6,0)}]
    \def\r{0.4775}
    \draw[very thick, bookaccent] (0,0) circle (\r);
    \fill[bookpitfall] (\r,0) circle (2pt) node[right, black] {$[0] = [1]$};
    \fill[bookdef] (0,\r) circle (2pt) node[above, black] {$[\tfrac14]$};
    \fill[bookexample] (-\r,0) circle (2pt) node[left, black] {$[\tfrac12]$};
    \draw[densely dashed, gray] (-\r,0) -- (\r,0);
  \end{scope}
\end{tikzpicture}
\caption{The curve $c$ wraps the real line around a circle of length $1$, drawn at the same scale as
the segment $[0, 1]$. All the integers land on the same point $[0]$. The points $[0]$ and $[\frac12]$ are at
distance $\frac12$ along the circle (half its length), but $\frac1\pi$ apart along the chord (dashed).}
\label{flat-circle:fig:circle}
\end{figure}

\section{Its length is \texorpdfstring{$1$}{1}}

Lengths on $\R/\Z$ are those of $\R$ transported by $\pi$. On any open interval $I$ of length less than $1$, $\pi$
is a bijection onto its image, so the infinitesimal length $\abs{dx}$ of $\R$ passes to $\R/\Z$: this is the
\emph{Riemannian metric} of $\R/\Z$. Walking once around, from $[0]$ to $[1] = [0]$, has length $1$. The
distance between two points is the length of the shorter way round,
\[
  d([x], [y]) = \min_{k\in\Z}\abs{y - x - k}\in\bigl[0, \tfrac12\bigr],
\]
the distance from $y - x$ to the nearest integer.

To see this circle in the plane with the right size, use
\[
  c(x) = \frac1{2\pi}\bigl(\cos 2\pi x, \sin 2\pi x\bigr),\qquad \norm{c'(x)} = \frac1{2\pi}\cdot2\pi = 1 .
\]
The curve $c$ has unit speed, so it preserves the length of every arc: it is an isometry of $\R/\Z$ onto
the circle of radius $\frac1{2\pi}$, whose perimeter is $1$. (The map $p$ of
\cref{flat-circle:prop:circle} has speed $2\pi$; it maps $\R/\Z$ onto the unit circle, of length $2\pi$, so it is
a homeomorphism but not an isometry.) Note that $d$ measures along the circle, not along the chord:
the points $[0]$ and $[\frac12]$ are at distance $\frac12$ in $\R/\Z$, but at distance $\frac1\pi$ in the plane.

\section{Curvature: extrinsic and intrinsic}

The circle $c(\R/\Z)$ is obviously curved in the plane: its unit tangent turns at the rate of $2\pi$ radians
per unit of length, the \emph{curvature} $\kappa = 1/r = 2\pi$ of a circle of radius $r$. This curvature is
\emph{extrinsic}: it describes how the curve sits in the ambient plane, and it changes if the same
curve is placed differently (deformed into an oval of the same length, its curvature varies from
point to point).

The \emph{intrinsic} geometry is what an inhabitant of the curve can measure using only lengths
along the curve. In dimension $1$ there is nothing to measure beyond arc length, and arc length
identifies a neighbourhood of every point isometrically with an interval of $\R$. Locally, $\R/\Z$ and $\R$
are indistinguishable: $\R/\Z$ is \emph{flat}. In the language of Riemannian geometry, the curvature tensor
measures how vectors fail to come back to themselves when transported around a small loop
spanned by two independent directions; in dimension $1$ there are no two independent directions,
and every one-dimensional Riemannian manifold is flat.

What the inhabitant \emph{can} detect is global: walking a distance $1$ in a straight line brings it back
to its starting point. This is a property of the topology (of the loops that cannot be shrunk to a
point), not of the curvature. ``Closing up on itself'' and ``being curved'' are different things: the
first is global and topological, the second local and metric. The bending we see in the picture
belongs to the drawing in $\R^2$, not to $\R/\Z$.

\section{The flat torus}

The same construction in dimension $2$ gives the \emph{flat torus} $\R^2/\Z^2$: the unit square with
opposite sides glued, as on the screen of the video games in which an object leaving on one side
comes back on the other (\cref{flat-circle:fig:torus}). It inherits the metric of $\R^2$, so it is intrinsically flat in the full
two-dimensional sense: small triangles have angles summing to $\pi$, and a neighbourhood of every point
is isometric to a piece of the plane. Here the intrinsic curvature could be nonzero, and it is zero.

\begin{figure}[ht]
\centering
\begin{tikzpicture}[>=Stealth, scale=2.4,
  mid1/.style={postaction={decorate}, decoration={markings, mark=at position 0.55 with {\arrow{>}}}},
  mid2/.style={postaction={decorate}, decoration={markings, mark=at position 0.5 with {\arrow{>>}}}}]
  \fill[bookaccent!6] (0,0) rectangle (1,1);
  \draw[thick, bookaccent, mid1] (0,0) -- (1,0);
  \draw[thick, bookaccent, mid1] (0,1) -- (1,1);
  \draw[thick, bookexample, mid2] (0,0) -- (0,1);
  \draw[thick, bookexample, mid2] (1,0) -- (1,1);
  % a straight line of slope 1/2 wrapping around
  \draw[very thick, bookdef] (0,0.2) -- (1,0.7);
  \draw[very thick, bookdef] (0,0.7) -- (0.6,1);
  \draw[very thick, bookdef] (0.6,0) -- (1,0.2);
  \fill (0,0.2) circle (0.6pt); \fill (1,0.2) circle (0.6pt);
  \node[below] at (0.5,0) {$x$};
  \node[left] at (0,0.5) {$y$};
\end{tikzpicture}
\caption{The flat torus $\R^2/\Z^2$: the unit square with the bottom side glued to the top one (single
arrows) and the left side to the right one (double arrows). A straight line (green, slope $\frac12$) leaving
through one side re-enters through the opposite side; it closes up after advancing by $2$
horizontally and $1$ vertically. Locally nothing distinguishes the square from the plane.}
\label{flat-circle:fig:torus}
\end{figure}

The familiar doughnut surface in $\R^3$ is a torus too, but with a different metric. Its Gaussian
curvature is positive on the outer part, negative on the inner part, and integrates to $0$
(Gauss--Bonnet: $\int K\,dA = 2\pi\chi = 0$ for a torus). The flat torus cannot be placed in $\R^3$ with its
lengths preserved by a twice differentiable surface, because every compact surface in $\R^3$ has a
point of positive Gaussian curvature (take the point farthest from the origin). It can be placed
isometrically in $\R^4$:
\[
  (x, y)\mapsto\frac1{2\pi}\bigl(\cos2\pi x, \sin2\pi x, \cos2\pi y, \sin2\pi y\bigr)
\]
is two copies of the curve $c$, one in each pair of coordinates. (A surprising theorem of Nash and
Kuiper even provides isometric embeddings into $\R^3$ that are only once differentiable, with
infinitely corrugated surfaces; they were first visualized by Borrelli, Jabrane, Lazarus and
Thibert in 2012.)

\section{Parameters that are angles}

Angles, phases, compass directions, the time of day or the day of the year are points of a circle,
and the rotation parameters of a model live on circles or tori. \Cref{gradient-tangent-space:exr:6}
showed what survives and what does not.
\begin{itemize}
  \item \emph{Gradient steps survive.} If $f$ is a differentiable function on $\R/\Z$, i.e.\ a $1$-periodic $\tilde f$ on
        $\R$, then $[x] - \eta\tilde f'(x) = [x - \eta\tilde f'(x)]$ is well defined. Locally the flat circle behaves like $\R$.
  \item \emph{Differences need a choice.} The displacement from $[x]$ to $[y]$ is defined only up to an
        integer. The usual choice is the representative in $[-\frac12, \frac12)$, of length $d([x], [y])$; it is
        the right one for small steps (momentum terms, finite differences) and wrong for large ones.
  \item \emph{Averages fail.} The arithmetic mean of representatives depends on the representatives:
        for the angles $0.95$ and $0.05$ (in turns), the mean of the representatives $0.95$ and $0.05$ is $0.5$,
        the opposite point, while the mean of $-0.05$ and $0.05$ is $0$. There is no affine structure, so
        there is no affine combination (\cref{gradient-tangent-space:prop:affine-comb}).
\end{itemize}
The standard remedy uses the embedding $c$. Points of the plane can be averaged, so average the
vectors $(\cos2\pi x_i, \sin2\pi x_i)$ and take the direction of the result:
\[
  \bar x = \frac1{2\pi}\operatorname{atan2}\Bigl(\sum_i\sin2\pi x_i,\ \sum_i\cos2\pi x_i\Bigr)\pmod 1 .
\]
For $0.95$ and $0.05$, the sines cancel and the cosines add up to $2\cos(0.1\pi)>0$, so $\bar x = 0$. This
\emph{circular mean} is undefined when the vectors cancel (for $0$ and $0.5$, which have no reasonable
mean). It is also the reason angles are usually given to a model as the pair $(\cos2\pi x, \sin2\pi x)$
rather than as a number in $[0, 1)$: the pair is continuous across $0\equiv1$, the number is not.
\end{notechapter}
